\usepackage{enumitem}
\usepackage{fontawesome5}

\tikzset{
	second box/.style={
			anchor=east,
			text=white,
			% rounded corners,
			fill=#1,
			xshift=-4mm,
		},
}

\tikzset{tcbreak/.style = {fill=white, thick, draw=#1, text=#1, font=\bfseries}}

\tcbset{
	common/.style n args={2}{
			colframe={#1},
			colback={#1!5},
			colbacktitle={#1},
			lower separated=false,
			coltitle=white,
			boxrule=1pt,
			fonttitle=\bfseries,
			enhanced,
			breakable,
			top=20pt,
			sharp corners,
			before skip=1em,
			after skip=2em,
			attach boxed title to top left={
					yshift=-0.25cm,
					xshift=0.38cm,
				},
			boxed title style={
					boxrule=0pt,
					colframe=white,
					arc=0pt,
					outer arc=0pt,
				},
			separator sign={~~},
			% -- overlays based on the following TeX SE answer: https://tex.stackexchange.com/a/545324/162087
			overlay unbroken={
					\node[text=white, align=right, fill=#1, xshift=-4mm, minimum height=6mm, anchor=east] at (frame.south east) {#2};
				},
			overlay first={
					\draw[thick, #1, decoration={coil, amplitude=0.5mm}, decorate]
					(frame.south west) -- (frame.south east);
					\path[font=\small\itshape] (frame.south) node [tcbreak={#1}] (cont) {continues in the next page \faHandPointRight};
				},
			overlay middle={%
					% upper break
					\draw[thick, #1, decoration={coil, amplitude=0.5mm}, decorate]
					(frame.north west) -- (frame.north east);
					\path[font=\small\itshape] (frame.north) node [tcbreak={#1}] (cont) {\faHandPointRight continues from the previous page};

					% lower break
					\draw[thick, #1, decoration={coil, amplitude=0.5mm}, decorate]
					(frame.south west) -- (frame.south east);
					\path[font=\small\itshape] (frame.south) node [tcbreak={#1}] (cont) {continues in the next page \faHandPointRight};
				},
			overlay last={
					\node[text=white, align=right, fill=#1, xshift=-4mm, minimum height=6mm, anchor=east] at (frame.south east) {#2};
					\draw[thick, #1, decoration={coil, amplitude=0.5mm}, decorate]
					(frame.north west) -- (frame.north east);
					\path[font=\small\itshape] (frame.north) node [tcbreak={#1}] (cont) {\faHandPointRight continues from the previous page};
				}
		},
	defstyle/.style={
			common={xpurple}{$\bm{\pi}$},
		},
	theoremstyle/.style={
			common={xgraycyan}{$\multimapdotbothA$},
		},
	lemmastyle/.style={
			common={xgrayblue}{$\multimap$},
		},
	proofstyle/.style={
			common={xgreen}{\textbf{QED}},
		},
	examplestyle/.style={
			common={xblue}{\faStar},
		},
	notestyle/.style={
			common={xred}{\textbf{!}},
		},
	challengestyle/.style={
			common={xgreen}{\textbf{?}},
		},
	quotestyle/.style={
			common={gray}{\textbf{"}},
		},
	summarystyle/.style={
			common={xpastelblue}{\textbf{\faClipboardList}},
		},
	motivatingproblemstyle/.style={
			common={xorange}{\textbf{\faRunning\faTrophy}},
		},
}

\newtcbtheorem[auto counter, number within=chapter]{definition}{Definition}{defstyle}{def}
\newtcbtheorem[auto counter, number within=chapter]{theorem}{Theorem}{theoremstyle}{theorem}
\newtcbtheorem[auto counter, number within=chapter]{lemma}{Lemma}{lemmastyle}{lemma}
\newtcbtheorem[auto counter, number within=chapter]{proof}{Proof}{proofstyle}{proof}
\newtcbtheorem[auto counter, number within=chapter]{example}{Example}{examplestyle}{example}
\newtcbtheorem[auto counter, number within=chapter]{note}{Note}{notestyle}{note}
\newtcbtheorem[auto counter, number within=chapter]{challenge}{Challenge}{challengestyle}{challenge}
\newtcbtheorem[auto counter, number within=chapter]{summary}{Summary}{summarystyle}{summary}
\newtcbtheorem[auto counter, number within=chapter]{motivation}{Motivating Problem}{motivatingproblemstyle}{motivation}
\newtcbtheorem[]{nquote}{Quote}{quotestyle}{quote}

% Bold text itemize/enumerate
% Taken from an answer to the TeX StackExchange question #266225
\newenvironment{descitemize}
{\begin{description}[leftmargin=*, before=\let\makelabel\descitemlabel]}
		{\end{description}}
\newcommand{\descitemlabel}[1]{\textbullet\ \textbf{#1}:}

% Box for other names side note
\tcbset{
	marginbox/.style={
		float,
		enhanced,
		size=fbox,
		colframe=#1!50,
		colback=#1!20,
		coltitle=white,
		title=synonyms,
	},
	evenmarginbox/.style={
		marginbox={#1}, 
		width=1in + \hoffset + \oddsidemargin,
	},
	oddmarginbox/.style={
		marginbox={#1}, 
		width=1in + \hoffset + \evensidemargin
	},
}
\NewDocumentCommand\synonyms{O{xpurple}mm}%
{
	% pass
}

% \newcommand{\synonyms}[3]{ % [1] is the term istself, [2] the list of synonyms
% 	\tcbox[
% 		enhanced,
% 		remember as=hltext,
% 		nobeforeafter,
% 		box align=base,
% 		boxrule=0pt,
% 		frame hidden,
% 		borderline west={1pt}{0pt}{#3},
% 		borderline south={1pt}{0pt}{#3},
% 		colback=#3!5,
% 		sharp corners,
% 		left=0pt,
% 		right=0pt,
% 		top=0pt,
% 		bottom=0pt,
% 	] {#1}
% 	\tcbox[
% 		float,
% 		floatplacement=t,
% 		enhanced,
% 		remember as=synonymsbox,
% 		boxrule=0pt,
% 		frame hidden,
% 		borderline west={8pt}{0pt}{#3},
% 		colback=#3!25,
% 		sharp corners,
% 		left=10pt,
% 		right=3pt,
% 		top=3pt,
% 		bottom=3pt,
% 	] {#2}
% 	\tikz[overlay, remember picture]{\draw[thick, #3!25] (hltext.north) -- (synonymsbox.north)}
% }
